\documentclass[11pt,a4paper]{report}
\usepackage[T1]{fontenc}
\usepackage[default]{sourcesanspro}


\usepackage[backend=biber,style=authoryear,sorting=nyt]{biblatex}
\addbibresource{references.bib}

\setcounter{secnumdepth}{3}
\setcounter{tocdepth}{3}

\title{Masters Dissertation}
\author{Nathan Calitz}
\date{\today}


\begin{document}
% Preliminary admin stuff
\maketitle
\pagenumbering{roman}
\chapter*{Plagarism and AI declaration}
\addcontentsline{toc}{chapter}{Plagiarism Declaration}

\chapter*{Abbreviations}
\addcontentsline{toc}{chapter}{Abbreviations}

\tableofcontents
\clearpage
\pagenumbering{arabic}

\chapter{Introduction}

Abstract:

Problem statement:

Aims, objectives, hypothesis:

Aims: To assess the current state of neurological disease drug discovery from plants and build a database combining this information

Objectives:
1. Assess some of the major natural product databases
2. Build a web accessible natural product database for neurological disease drug discovery
3. Test the data on wand to validate it and predict possible medications

Scope/outline dissertation

\chapter{Literature review}
\subsection{neurological disease}
Neurological diseases encompass a wide variety of ailments affecting the central and peripheral nervous systems. This includes neuroimmunological diseases such as multiple sclerosis and neuromyelitis spectrum disorder, degenerative diseases like dementia, Alzheimer’s, amyotrophic lateral sclerosis, and Parkinson’s, as well as psychiatric disorders like schizophrenia, anxiety, and bipolar disorder (Ropper et al., 2023). For the purposes of this work neurological disease is defined as all disorders with ICD-11 codes starting with a 6 (mental, behavioral or neurodevelopmental disorders) or an 8 (diseases of the nervous system) (The Lancet, 2018). Despite separation in the ICD-11 acknowledges the limitations of separating the two(pg 28 or 46) (World Health Organization, 2024) the describing psychological disorders as neurological disorders is inline with biological understanding (White et al., 2012) and has been done in other work (Luo et al., 2024).
Neurological diseases are the leading cause of disability worldwide and affect 3.4 billion people (Steinmetz et al., 2024).

\chapter{Databases in pytochemical drug discovery}
\chapter{Wand development}
\chapter{Drug discovery}
\chapter{Conclusion}



\end{document}